\subsection*{Despliegue en la nube}
\addcontentsline{toc}{subsection}{Despliegue en la nube}

\us[
	identificador=US-DE-01,
	titulo={Despliegue en la nube del front-end en el entorno de desarollo},
	como={Desarrollador},
	quiero={Poder acceder a la aplicación desde cualquier dispositivo con acceso a internet},
	requerimientos={\reqref{RNF-DE-01}, \reqref{RNF-DE-02}, \reqref{RNF-DE-05}, \reqref{RNF-SE-03}, \reqref{RNF-SE-05}},
	puntos={8},
	criterios={
			\begin{itemize}
				\item Existe un entorno de desarrollo en la nube.
				\item Cada nuevo commit a la rama main se despliega automáticamente a la nube.
				\item La aplicación es accesible a través de una URL de Internet.
				\item Sólo el equipo de desarrollo tiene acceso al entorno de desarrollo en la nube.
				\item Las credenciales de despliegue en la nube se almacenan de forma segura y no se exponen en el código fuente.
				\item La aplicación de desarrollo se sirve desde un subdominio del dominio principal del proyecto.
			\end{itemize}
		},
	tareas={
			\begin{itemize}
				\item Selección de un proveedor de nube para el entorno de desarrollo.
				\item Diseño de la infraestructura en la nube necesaria para el entorno.
				\item Configuración de la infraestructura como código.
				\item Creación de la infraestructura en el ciclo de CD (Continuous Deployment).
				\item Despliegue del código fuente del front-end a la nube.
				\item Configuración de la autenticación para el acceso al entorno de desarrollo en la
				      nube.
				\item Configuración del dominio personalizado para el entorno de desarrollo.
			\end{itemize}
		}
]

\us[
	identificador=US-DE-02,
	titulo={Despliegue en la nube del front-end en el entorno de producción},
	como={Usuario},
	quiero={Poder acceder a la aplicación desde cualquier dispositivo con acceso a internet},
	requerimientos={\reqref{RNF-CO-01}, \reqref{RNF-DE-01}, \reqref{RNF-DE-02}, \reqref{RNF-DE-03}, \reqref{RNF-DE-05}, \reqref{RNF-RE-04}},
	puntos={8},
	criterios={
			\begin{itemize}
				\item Existe un entorno de producción en la nube.
				\item El despliegue en producción se realiza automáticamente, pero se lanza manualmente.
				\item La aplicación es accesible a través de una URL de Internet.
				\item La infraestructura de la aplicación en producción es idéntica a la del entorno
				      de desarrollo.
				\item La aplicación de producción se sirve desde el dominio principal del proyecto.
			\end{itemize}
		},
	tareas={
			\begin{itemize}
				\item Creación de la infraestructura de producción a partir de la infraestructura de
				      desarrollo.
				\item Configuración de despliegue automático para producción.
				\item Configuración del dominio personalizado para el entorno de producción.
			\end{itemize}
		}
]

\us[
	identificador=US-DE-03,
	titulo={Despliegue en la nube de la API},
	como={Desarrollador},
	quiero={Poder acceder a los datos almacenados desde cualquiera de los entornos de la aplicación},
	requerimientos={\reqref{RNF-CO-01}, \reqref{RNF-DE-01}, \reqref{RNF-DE-02}, \reqref{RNF-DE-03}, \reqref{RNF-DE-05}, \reqref{RNF-SE-06}},
	puntos={5},
	criterios={
			\begin{itemize}
				\item Existe infraestructura en la nube para la API.
				\item Cada nuevo commit a la rama main se despliega automáticamente a la nube.
				\item La API es accesible a través de una URL de Internet.
				\item La API sólo admite peticiones desde el front-end de su propio entorno.
			\end{itemize}
		},
	tareas={
			\begin{itemize}
				\item Diseño de la infraestructura en la nube necesaria para la API.
				\item Configuración de la infraestructura como código.
				\item Creación de la infraestructura en el ciclo de CD (Continuous Deployment).
				\item Despliegue del código fuente de la API a la nube.
				\item Restricción del acceso a la API mediante el control de acceso al origen de
				      CloudFront, que la hace alcanzable únicamente a través del front-end de su
				      propio entorno.
			\end{itemize}
		}
]

\us[
	identificador=US-DE-04,
	titulo={Monitorización de la infraestructura en la nube},
	como={Desarrollador},
	quiero={Conocer desde fuera qué hace el sistema desplegado: qué peticiones atiende la API, cuánto tarda en atenderlas y qué errores se producen mientras se ejecuta},
	requerimientos={\reqref{RNF-MO-01}, \reqref{RNF-MO-03}, \reqref{RNF-MO-04}, \reqref{RNF-DE-03}},
	puntos={5},
	criterios={
			\begin{itemize}
				\item Los logs de las funciones Lambda se recogen en CloudWatch con un periodo de retención definido, en los dos entornos.
				\item La función de la API se ejecuta con el rastreo de X-Ray activo, de modo que sus trazas distinguen la inicialización del entorno de ejecución del tiempo consumido por el código.
				\item El entorno de producción cuenta con un panel de CloudWatch que reúne las invocaciones, los errores y la latencia de la función Lambda junto con las peticiones y las respuestas de error de la distribución de CloudFront.
				\item Los recursos de monitorización se declaran como infraestructura como código, y el entorno que dispone de panel se determina mediante un parámetro de la plantilla.
				\item Existe documentación sobre el acceso al panel, a los logs y a las trazas.
			\end{itemize}
		},
	tareas={
			\begin{itemize}
				\item Definición del periodo de retención de los grupos de logs existentes.
				\item Selección de las métricas que describen el estado de la API y de la distribución.
				\item Activación del rastreo de X-Ray en la función de la API, con los permisos que su envío de segmentos requiere.
				\item Creación del panel de CloudWatch en la plantilla de CloudFormation.
				\item Documentación del acceso al panel, a los logs y a las trazas.
			\end{itemize}
		}
]

\us[
	identificador=US-DE-05,
	titulo={Instrumentación de la API},
	como={Desarrollador},
	quiero={Saber en qué operación de la API se invierte el tiempo de cada petición y en cuál se produce un fallo, más allá de la duración total de la invocación},
	requerimientos={\reqref{RNF-MO-03}, \reqref{RNF-MO-01}},
	puntos={3},
	criterios={
			\begin{itemize}
				\item Las trazas de X-Ray de la función de la API incluyen un \textit{span} por petición, con el nombre de la ruta que la atiende.
				\item Cada petición se desglosa en la resolución de dependencias, la ejecución de la operación y la serialización de la respuesta.
				\item Una excepción no controlada queda registrada en el \textit{span} de la operación en la que se produce.
				\item La instrumentación sigue la decisión de muestreo de Lambda y no incrementa el número de trazas registradas.
				\item La aplicación se ejecuta en local sin cambios y sin enviar telemetría.
				\item Existe documentación sobre la configuración de la instrumentación y su coste.
			\end{itemize}
		},
	tareas={
			\begin{itemize}
				\item Evaluación del SDK de X-Ray frente a OpenTelemetry, a la vista del fin de soporte del primero.
				\item Actualización de FastAPI a la versión que emite \textit{spans} de OpenTelemetry de forma nativa.
				\item Incorporación del \textit{layer} de OpenTelemetry de AWS a la función de la API en la plantilla de CloudFormation, con las variables de entorno que lo configuran.
				\item Exclusión de \texttt{opentelemetry-api} del paquete de despliegue, para que la versión que se ejecute sea la del \textit{layer}.
				\item Ampliación de los permisos del rol de despliegue para asociar el \textit{layer} a la función.
				\item Documentación de la instrumentación, de su funcionamiento y de su efecto sobre el arranque en frío.
			\end{itemize}
		}
]
